\textbf{结论名称：} 关于直线 $y=x$ 的对称变换与反函数性质

\textbf{适用条件：} 点与曲线在同一平面直角坐标系中；反函数存在要求每个函数值只对应唯一原像；自反条件在全部定义域成立且复合有意义。

\textbf{变量说明：} $P(a,b)$ 为任意点，$D$ 为原函数定义域，$E=f(D)$ 为值域，$f^{-1}$ 为反函数。

\textbf{核心公式：}
\[
\boxed{(a,b)\longleftrightarrow(b,a)},\qquad
\boxed{F(x,y)=0\longleftrightarrow F(y,x)=0}.
\]
\[
 y=f(x)\longleftrightarrow x=f(y),\qquad
 f^{-1}(f(x))=x\quad(x\in D).
\]
\[
 \boxed{f(f(x))=x\ (x\in D)}
 \quad\Longleftrightarrow\quad
 \text{原函数图象自身关于 $y=x$ 对称}.
\]
上式要求 $f(D)\subseteq D$；成立后自动得到 $f(D)=D$。

\textbf{使用场景：} 求镜像方程、反函数、抽象函数求值、交点、切线斜率及圆锥曲线变换。

\textbf{简单例子：} $y=2^x$ 与 $y=\log_2 x$ 互为镜像；$y=1/x$ 在 $x\ne0$ 上自身对称。
